
\section{每个\trSmooth{}三次曲面（$\bar{k}$ 上）都是胀开平面}

我们现在来证明\mbox{定理~27.0.2}。正如 \S27.0.3 中所述，本节与上一节出人意料地相互独立；我们唯一需要的只是\mbox{练习~27.2.J}，而它也有可能用别的方式来证明。

设 $X$ 是（$\bar{k}$ 上的）一张\trSmooth{}三次曲面。设 $\ell$ 是 $X$ 上的一条直线，并在周围的空间 $\mathbb{P}^3$ 上选取\trCoordinate{}，使 $\ell$ 由 $x_0$ 与 $x_1$ 切出。从 $\ell$ 出发的\trProjection{}给出一个\trRationalMap{} $\mathbb{P}^3\dashrightarrow\mathbb{P}^1$（由 $[x_0,x_1,x_2,x_3]\mapsto[x_0,x_1]$ 给出），它延拓为 $X$ 上的一个\trMorphism{}。缘由在于：这个\trRationalMap{}可以通过胀开\trClosedSubscheme{} $V(x_0,x_1)$ 而得到解消（\mbox{练习~22.4.L}）。但 $(x_0,x_1)$ 在 $X$ 上切出的正是\trEffectiveCartierDivisor{} $\ell$，而胀开一个\trEffectiveCartierDivisor{}并不改变 $X$（\mbox{观察~22.2.1}）。

现在按\mbox{练习~27.2.J} 那样取定两条不交的直线 $\ell$ 与 $\ell'$，并考虑\trMorphism{} $\rho\colon X\to\mathbb{P}^1\times\mathbb{P}^1$，其中到第一个 $\mathbb{P}^1$ 的映射是从 $\ell$ 出发的\trProjection{}，而到第二个 $\mathbb{P}^1$ 的映射是从 $\ell'$ 出发的\trProjection{}。于是，第一个 $\mathbb{P}^1$ 可以与 $\ell'$ 等同起来，而第二个可以与 $\ell$ 等同起来。

\textbf{27.3.A.（练习）}\quad 证明：\trMorphism{} $\rho$ 是\trBirational{}的。提示：给定一般的点 $(p,q)\in\ell'\times\ell$\footnote{原书此处作「Given a general point of $(p,q)\in\ell'\times\ell$」，「of」当为衍文，径删。——译者注}，我们如下得到 $X$ 的一个点。$\mathbb{P}^3$ 中的直线 $pq$ 与三次曲面 $X$ 相交于三个点（由\mbox{\trBezout{}定理~9.3.D}）：$p$、$q$，以及某个第三点 $x\in X$；把 $(p,q)$ 送到 $x$。（这个想法早先在三次\trCurve{}的\trGroupLaw{}的展开中出现过；见\mbox{命题~19.10.3}。）反过来，给定一般的点 $x\in X$，从 $\ell'$ 与 $\ell$ 出发作\trProjection{}，便得到一个点 $(p,q)\in\ell'\times\ell$。

特别地，我们由此首次证明了：$\bar{k}$ 上每个\trSmooth{}三次曲面都是\trRational{}的。

\textbf{27.3.B.（练习：$\rho$ 恰好收缩 $\ell_1,\dots,\ell_5$）}\quad 证明：$\rho$ 在\mbox{练习~27.2.J} 所提到的诸 $\ell_i$ 之外是\trIsomorphism{}，且每个 $\rho(\ell_i)$ 都是一个点 $p_i\in\ell'\times\ell$。

\textbf{27.3.1. \trhProposition{}。}\ ——\ \emph{\trMorphism{} $\rho\colon X\to\mathbb{P}^1\times\mathbb{P}^1$ 是 $\mathbb{P}^1\times\mathbb{P}^1$ 在这五个 $p_i$ 处的\trBlowup{}。}

我们给出两个证明。第一个展示了若干技术工具——它们是我们已经发展出来的（或将来才会发展的——此处有一处「前向引用」）——的效用。第二个则是脚踏实地的显式「计算」。不同的人偏爱不同的路径，但把二者并排对照着看自有其好处。第一个表明：在我们对几何并无显式掌控的相当一般的情形下，如何运用这些工具；第二个则让我们对所论的几何有相当完全而直接的掌控，且无须用到形式函数。

\textbf{27.3.2. \mbox{命题~27.3.1} 的证明一。}\quad 由\trCastelnuovo{}判别准则（\mbox{定理~28.6.1}），诸直线 $\ell_i$ 都是 $(-1)$-曲线（\mbox{练习~27.1.B}），故它们可以被收缩。更精确地说：存在\trMorphism{} $\beta\colon X\to X'$，它是 $X'$ 在五个\trClosedPoint{} $p_1',\dots,p_5'$ 处的\trBlowup{}，使得 $\ell_i$ 正是 $p_i'$ 之上的\trExceptionalDivisor{}。我们希望证明 $X'$ 就是 $\mathbb{P}^1\times\mathbb{P}^1$。\trMorphism{} $\rho\colon X\to\mathbb{P}^1\times\mathbb{P}^1$ 给出\trMorphism{} $\rho'\colon X'\setminus\{p_1',\dots,p_5'\}\to\mathbb{P}^1\times\mathbb{P}^1$\footnote{原书此处作「$X'\setminus\{p_1',\ldots\,p_5'\}$」，脱一逗号，当作「$p_1',\ldots,p_5'$」，径补。——译者注}。我们现在证明：$\rho'$ 对每个 $i$ 都能延拓过 $p_i'$，并把 $p_i'$ 送到 $p_i$。选取 $p_i\in\mathbb{P}^1\times\mathbb{P}^1$ 的一个与 $\mathbb{A}^2$ \trIsomorphic{}的开邻域，取\trCoordinate{} $x$ 与 $y$。此时 $x$ 与 $y$ 都拉回为 $p_i'$ 的某个刺穿的开邻域上的函数（也就是说，存在 $p_i'$ 的某个开邻域 $U$，使得 $x$ 与 $y$ 是 $U\setminus\{p_i'\}$ 上的函数）。由\mbox{代数\trHartogs{}引理~13.5.19}，它们可以延拓过 $p_i'$；又因 $\mathbb{P}^1\times\mathbb{P}^1$ 是\trSeparated{}的，这个延拓是唯一的——如果你实在需要，可以动用\mbox{「既约到分离」定理~11.3.2}。于是 $\rho'$ 延拓过 $p_i'$。（你看出为什么 $\rho'(p_i')=p_i$ 吗？）

\textbf{27.3.C.（练习）}\quad 证明：\trBirational{}\trMorphism{} $\rho'\colon X'\to\mathbb{P}^1\times\mathbb{P}^1$ 是可逆的。提示：你可以使用\trZariskiMainTheorem{}（以\mbox{练习~28.5.E} 的形式），但无须动用如此强力的工具。不如这样：注意\trBirational{}映射 $\rho'^{-1}$ 在 $p_1,\dots,p_5$ 之外是\trMorphism{}，再用与上一段本质相同的论证，把 $\rho'^{-1}$ 延拓过每个 $p_i$。\hfill$\square$

\textbf{27.3.3. \mbox{命题~27.3.1} 的证明二。}\quad （这个非常经典的论证的这一利落版本，是我从 J.~Koll\'ar 那里学来的；Koll\'ar 以善于用简单的方式解释复杂的事物而闻名。）在 $\mathbb{P}^3$ 上选取\trCoordinate{}，使 $\ell$ 由 $x_0=x_1=0$ 给出，而 $\ell'$ 由 $x_2=x_3=0$ 给出。令 $\ell$ 上的射影\trCoordinate{}为 $u_0$ 与 $u_1$，$\ell'$ 上的射影\trCoordinate{}为 $v_0$ 与 $v_1$（于是，比如说，\trInclusionMap{} $\ell\hookrightarrow\mathbb{P}^3$ 由 $[u_0,u_1]\mapsto[0,0,u_0,u_1]$ 给出）。考虑\trRationalMap
\[
\alpha\colon\mathbb{P}^1_{u_0,u_1}\times\mathbb{P}^1_{v_0,v_1}\times\mathbb{P}^1_{s,t}\longrightarrow\mathbb{P}^3
\]
（其中诸下标标明射影\trCoordinate{}的名字），它由
\[
\tag{27.3.3.1}\label{eq:alpha-cubic-map}
([u_0,u_1],\,[v_0,v_1],\,[s,t])\mapsto[sv_0,\,sv_1,\,tu_0,\,tu_1].
\]
给出。

由于 $X$ 包含 $\ell=V(x_0,x_1)$ 与 $\ell'=V(x_2,x_3)$，故 $X\subset\mathbb{P}^3$ 的方程形如
\[
F(x_0,x_1,x_2,x_3)=L_{02}x_0x_2+L_{03}x_0x_3+L_{12}x_1x_2+L_{13}x_1x_3,
\]
其中诸 $L_{ij}$ 是 $x_0$、$x_1$、$x_2$、$x_3$ 的线性型。于是
\[
\alpha^*F=st\,\bigl(sG(u_0,u_1,v_0,v_1)+tH(u_0,u_1,v_0,v_1)\bigr).
\]
由 \eqref{eq:alpha-cubic-map}，$G$ 对诸 $u_i$ 是 $1$ 次的、对诸 $v_i$ 是 $2$ 次的，而 $H$ 对诸 $u_i$ 是 $2$ 次的、对诸 $v_i$ 是 $1$ 次的。

\textbf{27.3.D.（练习）}\quad 把\trMorphism{} $X\to\mathbb{P}^1\times\mathbb{P}^1$ 解读为一个\trRationalMap{}：先用 $\alpha$ 拉回 $F=0$，再投射到 $\mathbb{P}^1_{u_0,u_1}\times\mathbb{P}^1_{v_0,v_1}$。证明：被 $X\to\mathbb{P}^1\times\mathbb{P}^1$ 收缩（到 $p_1,\dots,p_5$）的那五条直线 $\ell_1,\dots,\ell_5$，在某种意义上必定（按集合论意义）对应于轨迹 $G=H=0$。证明：$G=H=0$ 按集合论意义切出的正是 $\mathbb{P}^1\times\mathbb{P}^1$ 中的五个点 $p_1,\dots,p_5$。证明：（在 $\mathbb{P}^1\times\mathbb{P}^1$ 的\trIntersectionTheory{}中，\mbox{练习~20.2.C}）$V(G)\cdot V(H)=5$。证明：在 $\mathbb{P}^1\times\mathbb{P}^1$ 的五点 $p_1,\dots,p_5$ 中任何一点附近，\trEffectiveCartierDivisor{} $V(G)$ 与 $V(H)$ 都必定横截地相交（具有「重数 $1$」）——你得把这件事说精确。\footnote{原书此处作「trasnversely」，当为「transversely」之误，径改。——译者注}证明：在每个 $p_i$ 附近，这个\trProjection{}必定是 $\mathbb{P}^1\times\mathbb{P}^1$ 在 $p_i$ 处的\trBlowup{}。\hfill$\square$

\textbf{27.3.4. \mbox{定理~27.0.2} 证明的收官。}\quad 作为推论，我们看到 $X$ 是 $\mathbb{P}^1\times\mathbb{P}^1$ 在五个点处的\trBlowup{}。由于 $\mathbb{P}^1\times\mathbb{P}^1$ 在一点处的\trBlowup{}\trIsomorphic{}于 $\mathbb{P}^2$ 在两点处的\trBlowup{}（\mbox{练习~22.4.K}），\mbox{定理~27.0.2} 随之成立。\hfill$\square$

\textbf{27.3.5. 把过程反过来。}

（这是 \mbox{22.4.10(iii)} 的一个更精确的版本。）这个过程可以反过来：我们可以在 $\mathbb{P}^2$ 的六个点处作\trBlowup{}，并把它嵌入 $\mathbb{P}^3$。我们先解释为什么不能在随随便便六个点处胀开 $\mathbb{P}^2$，就指望把结果嵌入 $\mathbb{P}^3$。由于这张三次曲面是反典范地嵌入的（\mbox{练习~27.1.A}），我们看到 $X$ 中的任何\trCurve{} $C$ 都必须满足 $\deg_C\mathscr{K}_X<0$。

\textbf{27.3.E.（练习）}\quad 设 $\mathbb{P}^2$ 依次在 $p_1,\dots,p_6$ 处被胀开，得到一张\trSmooth{}\trSurface{} $X$。
（a）若 $p_i$ 落在胀开 $p_j$ 所得的\trExceptionalDivisor{}上（$i>j$），证明：存在\trIsomorphic{}于 $\mathbb{P}^1$ 的\trCurve{} $C\subset X$，使得 $\deg_C\mathscr{K}_X\ge 0$。
（b）若诸 $p_i$ 是 $\mathbb{P}^2$ 中互不相同的点，而其中有三点共线，证明：存在\trIsomorphic{}于 $\mathbb{P}^1$ 的\trCurve{} $C\subset X$，使得 $\deg_C\mathscr{K}_X\ge 0$。
（c）若这六个 $p_i$ 是一条\trSmooth{}\trConic{}上互不相同的点，证明：存在\trIsomorphic{}于 $\mathbb{P}^1$ 的\trCurve{} $C\subset X$，使得 $\deg_C\mathscr{K}_X\ge 0$。

于是，要想通过在 $\mathbb{P}^2$ 上胀开六个点来得到一张\trSmooth{}三次曲面，我们唯一的希望就是：胀开六个互不相同的点，其中任何三点不共线，且它们不全在一条\trConic{}上。

\textbf{27.3.6. \trhProposition{}。}\ ——\ \emph{在六个互不相同的点处胀开的 $\mathbb{P}^2$——其中任何三点不共线，且它们不全在一条\trConic{}上——其反典范映射给出一个到 $\mathbb{P}^3$ 的\trClosedEmbedding{}，其像是一张三次曲面。}

由于我们并不会用到这条命题，这里只描述其证明的主要步骤：首先，数一数反典范丛 $\mathscr{K}_{\mathbb{P}^2}^{\vee}\cong\OO_{\mathbb{P}^2}(3)$ 的\trSection{}（$\mathbb{P}^2$ 上在这六个点处消失的三次型构成一个四维\trVectorSpace{}，而经由\mbox{练习~22.4.S(a)}，它们相应于胀开曲面的反典范丛的\trSection{}）。然后证明这些\trSection{}分离 $X$ 的点与切向量，从而证明反典范\trLinearSystem{}给出一个\trClosedEmbedding{}，即\mbox{定理~19.1.1}。在这一步中，明智地使用\trCremona{}变换（\mbox{练习~7.5.J}）可以减少繁琐的分情况讨论。

\textbf{27.3.F.（练习）}\quad 设 $X$ 是 $\mathbb{P}^2_{\bar{k}}$ 在六个互不相同的点 $p_1,\dots,p_6$ 处的\trBlowup{}，其中任何三点不共线，且它们不全在一条\trConic{}上。验证：$X$ 上仅有的 $(-1)$-曲线是六个\trExceptionalDivisor{}、15 条直线 $p_ip_j$ 的\trProperTransform{}，以及过这六点中任意五点的那六条\trConic{}的\trProperTransform{}，总计 $27$ 条。

\textbf{27.3.G.（练习）}\quad 再把\mbox{练习~27.2.I} 与\mbox{练习~27.2.J} 解一遍，这一次利用 $X$ 作为 $\mathbb{P}^2$ 之\trBlowup{}的描述。

\textbf{27.3.7.}\ \emph{\trhRemark{}。}\quad 若你在 $\mathbb{P}^2$ 上胀开 $4\le n\le 8$ 个点，其中任何三点不共线、任何六点不在一条\trConic{}上，那么诸直线之构型的对称群是一个\trWeyl{}群，如下表所示。

\begin{center}
\begin{tabular}{|c||c|c|c|c|c|}
\hline
$n$ & $4$ & $5$ & $6$ & $7$ & $8$ \\\hline
 & $W(A_4)$ & $W(D_5)$ & $W(E_6)$ & $W(E_7)$ & $W(E_8)$ \\\hline
\end{tabular}
\end{center}

（如果你了解\trDynkinDiagram{}，你或许能看出其中的模式，并或许能诠释 $n=3$ 与 $n=9$ 时会发生什么。相应于 $n=4$ 的那张曲面出现在 \S19.8.3 中。）这推广了\mbox{注记~27.2.5} 的一部分，其余部分也可以类似地推广。

\textbf{27.3.H.（练习）}\quad 证明：\mbox{注记~27.3.7} 中所描述的那些曲面都是\trDelPezzo{}曲面（在\mbox{练习~21.5.H} 中定义）。（做这件事无须了解\trDynkinDiagram\footnote{原书此处作「You don't know need to know about Dynkin diagrams」，「know」当为衍文，径删。——译者注}。事实证明，每个\trDelPezzo{}曲面都属于这里所描述的、或\mbox{练习~21.5.H} 中所描述的类型。）
